% !TEX root = ../main.tex
\chapter[Testes e validação]{Testes e validação do sistema}
\label{chap:testes}

\section{Objetivo e escopo dos testes}

Os testes de ponta a ponta foram desenvolvidos com Cypress para verificar os principais fluxos do FlashFood pela interface da aplicação. Os cenários abrangem cadastro e autenticação de usuários, controle de acesso, catálogo de produtos, carrinho de compras, pedidos, histórico de mesas e entrega. Os arquivos estão organizados no diretório \texttt{cypress/e2e/}.

\section{Arquivos e cenários de teste}

As tabelas a seguir apresentam os arquivos Cypress e a finalidade de cada teste. A coluna de requisitos registra apenas as associações indicadas para os cenários; o símbolo -- indica que não foi informado um requisito específico.

\subsection{Cadastro, autenticação e controle de acesso}

\begin{table}[htbp]
\centering
\caption{Testes de cadastro, autenticação e controle de acesso}
\label{tab:testes-acesso}
\footnotesize
\begin{tabularx}{\textwidth}{c p{3.7cm} p{1.5cm} c X}
\hline
\textbf{ID} & \textbf{Arquivo} & \textbf{Requisito} & \textbf{Aprovados} & \textbf{O que o teste faz} \\
\hline
T01 & \texttt{AcessoAdmin.cy.js} & RF03 & 1 & Tenta abrir o painel administrativo sem autenticação e verifica se o sistema redireciona o visitante para a tela de login. \\
\hline
T02 & \texttt{AcessoPerfil.cy.js} & RF03 & 2 & Verifica que um cliente não consegue abrir o gerenciamento de produtos e que um administrador autorizado consegue acessar essa funcionalidade. \\
\hline
T03 & \texttt{AcessoProtegido.cy.js} & RF03 & 8 & Acessa oito páginas protegidas, incluindo perfil, pedidos, carrinho, mesas e entregas, sem uma sessão ativa e verifica o redirecionamento para o login. \\
\hline
T04 & \texttt{Cadastro.cy.js} & RF01 & 1 & Cadastra um cliente pela página de registro e, em seguida, confirma que ele consegue entrar no sistema. \\
\hline
T05 & \texttt{CriarFuncionario.cy.js} & RF01 & 1 & Verifica o cadastro de um funcionário por meio do painel administrativo. \\
\hline
T06 & \texttt{DeletarFuncionario.cy.js} & -- & 1 & Verifica a remoção de um funcionário pelo painel administrativo. \\
\hline
T07 & \shortstack[l]{\texttt{LogandoComFuncionario}\\\texttt{Criado.cy.js}} & RF01, RF02 & 2 & Cadastra um funcionário autorizado e verifica se ele consegue se autenticar no sistema. \\
\hline
T08 & \texttt{Logout.cy.js} & RF02 & 1 & Encerra a sessão pelo menu e verifica se uma página protegida volta a exigir autenticação. \\
\hline
T09 & \texttt{ValidacaoFormularios.cy.js} & -- & 4 & Verifica campos obrigatórios nas telas de login, cadastro e redefinição de senha, além do link para voltar ao login. \\
\hline
T10 & \texttt{loginAdmin.cy.js} & RF02 & 1 & Autentica um administrador e verifica se o painel administrativo é aberto. \\
\hline
T11 & \texttt{LoginInvalido.cy.js} & RF02 & 1 & Tenta entrar com credenciais inválidas e verifica se o sistema permanece na tela de login. \\
\hline
\end{tabularx}
\legend{Elaborado pelo autor (2026).}
\end{table}

\subsection{Produtos, cardápio, carrinho e pedidos}

\begin{table}[htbp]
\centering
\caption{Testes de produtos, cardápio, carrinho e pedidos}
\label{tab:testes-produtos-pedidos}
\footnotesize
\begin{tabularx}{\textwidth}{c p{3.7cm} p{1.5cm} c X}
\hline
\textbf{ID} & \textbf{Arquivo} & \textbf{Requisito} & \textbf{Aprovados} & \textbf{O que o teste faz} \\
\hline
T12 & \texttt{AtualizarProduto.cy.js} & -- & 1 & Abre a edição de um produto no painel administrativo, altera seus dados e verifica se a atualização foi realizada. \\
\hline
T13 & \texttt{BuscaProduto.cy.js} & -- & 1 & Pesquisa produtos pelo nome na tabela administrativa e verifica se os resultados são filtrados. \\
\hline
T14 & \texttt{Carrinho.cy.js} & RF06 & 1 & Adiciona um produto ao carrinho, altera a quantidade, confere o subtotal e remove o produto. \\
\hline
T15 & \texttt{CatalogoProdutos.cy.js} & RF05 & 1 & Cadastra um produto disponível e verifica se ele aparece no catálogo de bebidas do cliente com preço e descrição. \\
\hline
T16 & \texttt{CriandoProduto.cy.js} & RF11 & 1 & Cria um produto pelo painel administrativo e verifica se seus dados aparecem na listagem. \\
\hline
T17 & \texttt{DeletandoProduto.cy.js} & -- & 1 & Verifica a remoção de um produto pelo painel administrativo. \\
\hline
T18 & \texttt{NavegacaoCardapio.cy.js} & RF05 & 7 & Verifica as páginas do cardápio, a navegação até a seção de bebidas e a exibição do controle do menu em uma tela pequena. \\
\hline
T19 & \texttt{PersonalizacaoPizza.cy.js} & -- & 3 & Verifica o formulário de preços por tamanho, o salvamento e a edição desses preços e a disponibilidade das bordas ativas para pizzas. \\
\hline
T20 & \texttt{FluxoCompraEntrega.cy.js} & RF04, RF07--RF11 & 1 & Percorre o fluxo de compra: registra o cliente e o pedido, confere os dados salvos, consulta e atualiza o pedido no painel administrativo, acompanha o status como cliente e verifica as etapas da entrega. \\
\hline
T21 & \texttt{MesasHistorico.cy.js} & -- & 2 & Verifica os resumos e filtros do histórico de mesas e confirma a aplicação do filtro por datas. \\
\hline
\end{tabularx}
\legend{Elaborado pelo autor (2026).}
\end{table}

\section{Ambiente e execução}

Os testes foram executados com Cypress, versão \texttt{15.18.0}, utilizando o navegador Google Chrome, versão \texttt{149}. O ambiente de execução utilizou Node.js, versão \texttt{24.18.0}. A aplicação foi executada em contêineres gerenciados pelo Docker Compose. O Cypress acessou o serviço \texttt{app} pelo endereço interno \texttt{http://app:8000}, configurado pela variável \texttt{CYPRESS\_BASE\_URL}.

\section{Resultados}

Na execução registrada, foram executados 21 arquivos de teste, totalizando 42 casos. Todos os 42 casos foram aprovados; não houve falhas, casos pendentes ou casos ignorados. A execução completa durou 1 minuto e 46 segundos.

\section{Evidências}

Como evidência da execução, foi considerado o relatório de resultados do Cypress enviado pelo autor, que registra 21 arquivos e 42 casos aprovados, sem falhas, pendências ou casos ignorados. Essa execução não gerou capturas de tela nem vídeo.
